% Generated from validated Article Draft Result. Do not hand-edit; revise the structured draft instead.

2022年の言語モデル史をparameter数だけで追うと、年内に独立したもう一つの設計軸を見落とす。InstructGPTは、pretraining済みモデルを人間の指示や選好へ合わせるpost-trainingを前景化した。一方、FLAN系のinstruction tuningは多数のtaskをinstruction形式で学習させる方向を押し広げ、Self-Instructはinstruction dataをモデル自身から生成・選別して拡張する方向を示した。ここで重要なのは、いずれも『もっと大きくpretrainする』こととは別の操作としてモデルの振る舞いを設計している点である。\autocite{src-0c551201e765,src-87477ad8b476,src-e58f1e9c8754,src-fadb2b92c538}


Chain-of-Thought Promptingが属するのは、さらに別の層である。これはモデルの重みを更新するpost-trainingそのものではなく、推論時のpromptによって中間的な推論過程を引き出す方法として提示された。したがって、InstructGPTとChain-of-Thoughtを一括して『reasoning modelの誕生』と読むより、weight-spaceで振る舞いを変える方法と、inference-timeに能力の出方を変える方法が並行して発達したと捉える方が、2022年の設計空間を正確に表す。\autocite{src-87477ad8b476,src-a17359c66f6b,src-e58f1e9c8754}


post-training側でも手段は一様ではない。InstructGPTが人間のdemonstrationやpreference signalを利用するRLHF系の流れを代表するのに対し、FLANはinstruction-finetuningの規模とtask多様性を拡張する方向、Self-Instructはsynthetic instruction dataを利用してデータ作成のボトルネックへ手を入れる方向として読める。これらは目的の一部を共有していても、教師信号、学習工程、依存するデータ生成過程が異なる。\autocite{src-0c551201e765,src-87477ad8b476,src-e58f1e9c8754,src-fadb2b92c538}


この年の変化は、能力を一つの方法で『追加する』話ではない。pretrainingで表現能力を作り、post-trainingで望ましい応答へ寄せ、promptingで推論時の挙動を変え、synthetic dataでinstruction coverageを広げるという複数のレバーが分離して見えるようになった。翌年以降の対話製品やagent的利用を理解するうえでも、2022年にこの層分けが可視化されたことが重要である。ただし、後年の成熟した製品形態を2022年の各研究へ逆投影してはならない。\autocite{src-0c551201e765,src-87477ad8b476,src-a17359c66f6b,src-e58f1e9c8754,src-fadb2b92c538}


\begin{claimboundary}
本節で扱う性能・有効性は各論文・各組織が一次資料で提示した範囲に帰属する。本サーベイは、RLHF、instruction tuning、prompting、synthetic instruction dataのどれが普遍的に優れているかを独立再現した順位へ換算しない。\autocite{src-87477ad8b476,src-a17359c66f6b,src-e58f1e9c8754}
\end{claimboundary}
